\section{Code, data and the Overleaf project}\label{app:links}
The manuscript source is in the Overleaf project linked \overleaflink. The code, the saved
result files and the notebook are at \url{https://github.com/utshabkg/scminer-clustering}.

\texttt{notebooks/assignment3\_results.ipynb} rebuilds every table and figure in this
paper from the saved result files and displays them with the rows behind them, so the
displays can be checked without rerunning an experiment. Appendix~\ref{app:repro} gives
the layout of those files and the commands that produced them.

\section{The datasets}\label{app:data}
Table~\ref{tab:data} lists the ten benchmark datasets with their sizes, label counts and
sparsity. Nine fall below the 5,000 cell point at which MICA switches from its small-data
mode to its graph mode, and PBMC14K is the only one above it, which is why mode and size
move together in the runs as shipped.

\begin{table}[t]
\centering
\caption{The ten benchmark datasets, as released by the scMINER authors and used here without further preprocessing. The smallest class column is that class as a share of all cells, and the last column is the share of the matrix that is zero.}
\label{tab:data}
\footnotesize
\setlength{\tabcolsep}{4pt}
\begin{tabular}{lrrrrr}
\toprule
Dataset & Cells & Genes & Types & Smallest (\%) & Zeros (\%) \\
\midrule
Yan & 124 & 20,286 & 8 & 2.4 & 38 \\
Goolam & 124 & 23,571 & 5 & 4.8 & 46 \\
Buettner & 182 & 8,989 & 3 & 31.9 & 38 \\
Pollen & 249 & 6,982 & 11 & 2.8 & 25 \\
Chung & 317 & 10,897 & 4 & 7.9 & 96 \\
Usoskin & 622 & 16,383 & 4 & 13.0 & 78 \\
Kolodziejczyk & 704 & 13,698 & 3 & 22.6 & 37 \\
Klein & 2,717 & 24,175 & 4 & 11.2 & 66 \\
Zeisel & 2,903 & 14,890 & 7 & 3.3 & 76 \\
PBMC14K & 13,605 & 8,846 & 7 & 13.5 & 94 \\
\bottomrule
\end{tabular}
\end{table}


\section{All four rules in one table}\label{app:rulesall}
Table~\ref{tab:rulesall} puts the two stability readings of Table~\ref{tab:rules} beside
the shipped defaults and the shared-space silhouette, so the spread across rules quoted in
Section~\ref{sec:results} can be read off directly.

\begin{table*}[t]
\centering
\caption{ARI under all four label-free rules, with the chosen cluster count in parentheses. The two stability columns repeat Table~\ref{tab:rules} and are included so the four rules can be compared in one place. The Default and Silhouette rules are deterministic given the input and were run once, so they carry no seed range. PBMC14K was run under the stability rule only for MICA, its shipped small-data default not applying in graph mode.}
\label{tab:rulesall}
\footnotesize
\setlength{\tabcolsep}{4pt}
\begin{tabular}{lcccccccc}
\toprule
 & \multicolumn{2}{c}{Default} & \multicolumn{2}{c}{Silhouette} & \multicolumn{2}{c}{Finest stable} & \multicolumn{2}{c}{Most stable} \\
\cmidrule(lr){2-3}\cmidrule(lr){4-5}\cmidrule(lr){6-7}\cmidrule(lr){8-9}
Dataset (types) & MICA & Leiden & MICA & Leiden & MICA & Leiden & MICA & Leiden \\
\midrule
Yan (8) & 0.765\,(4) & 0.845\,(5) & 0.794\,(5) & 0.765\,(4) & 0.496\,(19\textrm{-}20) & 0.845\,(5\textrm{-}8) & 0.496\,(3) & 0.765\,(4\textrm{-}5) \\
Goolam (5) & 0.582\,(4) & 0.349\,(6) & 0.630\,(2) & 0.619\,(2) & 0.593\,(6\textrm{-}7) & 0.429\,(2\textrm{-}3) & 0.630\,(2\textrm{-}3) & 0.619\,(2\textrm{-}3) \\
Buettner (3) & 0.919\,(4) & 0.237\,(4) & 0.474\,(2) & 0.154\,(2) & 0.934\,(3) & 0.232\,(3) & 0.474\,(2\textrm{-}3) & 0.134\,(2\textrm{-}3) \\
Pollen (11) & 0.493\,(4) & 0.865\,(9) & 0.942\,(12) & 0.870\,(9) & 0.898\,(15\textrm{-}18) & 0.755\,(7\textrm{-}11) & 0.087\,(2) & 0.755\,(5\textrm{-}11) \\
Chung (4) & 0.655\,(4) & 0.567\,(7) & 0.551\,(9) & 0.046\,(2) & 0.474\,(11\textrm{-}12) & 0.129\,(3\textrm{-}4) & 0.100\,(2) & 0.129\,(3) \\
Usoskin (4) & 0.839\,(4) & 0.393\,(10\textrm{-}11) & 0.347\,(2) & 0.761\,(3) & 0.818\,(3) & 0.633\,(6) & 0.818\,(3) & 0.757\,(3\textrm{-}6) \\
Kolodziejczyk (3) & 0.834\,(4) & 0.525\,(7) & 0.727\,(3) & 0.525\,(7) & 0.722\,(3\textrm{-}5) & 0.838\,(4) & 0.630\,(2) & 0.838\,(4) \\
Klein (4) & 0.980\,(4) & 0.469\,(10\textrm{-}11) & 0.501\,(2) & 0.859\,(3) & 0.501\,(2\textrm{-}3) & 0.973\,(4) & 0.501\,(2) & 0.859\,(3\textrm{-}4) \\
Zeisel (7) & 0.478\,(4) & 0.378\,(21) & 0.386\,(2) & 0.473\,(3) & 0.518\,(3) & 0.855\,(8\textrm{-}9) & 0.518\,(2\textrm{-}3) & 0.855\,(8) \\
PBMC14K (7) & n/a & 0.866\,(8) & n/a & n/a & 0.629\,(5) & 0.397\,(4) & 0.629\,(5) & 0.397\,(4) \\
\midrule
\textit{Mean, nine small} & \textbf{0.727} & \textbf{0.514} & \textbf{0.595} & \textbf{0.564} & \textbf{0.661} & \textbf{0.632} & \textbf{0.473} & \textbf{0.635} \\
\bottomrule
\end{tabular}
\end{table*}


\section{The stability rule on a common range}\label{app:common}
The 2 to 20 filter caps both grids without matching them. Table~\ref{tab:common} gives, for each
dataset, the span of cluster counts the baseline's resolution grid produced, and the stability
rule re-applied to MICA's saved scores restricted to the set of counts the baseline produced in
the same run. That set can have gaps: on Yan the baseline makes 2 and 4 to 8 clusters but never
3. Nothing is rerun and no labels are used to restrict, but the restriction was chosen after the
results were seen and is applied to MICA only, which is why the matched runs of
Section~\ref{sec:matched} replace it as the comparison we rely on. The largest gains in MICA's finest-stable ARI are
on Yan, where its pick falls from 19 clusters to 7 and its ARI rises from 0.496 to 0.925, and on
Klein, by 0.36. On Pollen the pick falls from 17 to the true 11 at 0.969, though 11 is also the
most clusters the baseline makes there.

\begin{table*}[t]
\centering
\caption{The stability rule on its own grids and on a common range. The second column is the span of cluster counts the baseline's resolution grid produced over the three seeds. MICA's $k$ covers 2 to 20 on every dataset. The middle pair re-applies the rule's own selection to MICA's saved scores restricted to the set of counts the baseline produced in the same run, which can have gaps within that span. The restriction uses no labels but was chosen after the results were seen and applies to MICA only. Values are medians over three seeds, F the finest-stable reading and M the plain maximum.}
\label{tab:common}
\footnotesize
\setlength{\tabcolsep}{4pt}
\begin{tabular}{lccccccc}
\toprule
 & & \multicolumn{2}{c}{MICA, own grid} & \multicolumn{2}{c}{MICA, common range} & \multicolumn{2}{c}{Leiden} \\
\cmidrule(lr){3-4}\cmidrule(lr){5-6}\cmidrule(lr){7-8}
Dataset (types) & Baseline span & F & M & F & M & F & M \\
\midrule
Yan (8) & 2\textrm{-}8 & 0.496 & 0.496 & 0.925 & 0.765 & 0.845 & 0.765 \\
Goolam (5) & 2\textrm{-}8 & 0.593 & 0.630 & 0.593 & 0.630 & 0.429 & 0.619 \\
Buettner (3) & 2\textrm{-}14 & 0.934 & 0.474 & 0.934 & 0.474 & 0.232 & 0.134 \\
Pollen (11) & 4\textrm{-}11 & 0.898 & 0.087 & 0.969 & 0.549 & 0.755 & 0.755 \\
Chung (4) & 2\textrm{-}11 & 0.474 & 0.100 & 0.477 & 0.512 & 0.129 & 0.129 \\
Usoskin (4) & 3\textrm{-}18 & 0.818 & 0.818 & 0.818 & 0.818 & 0.633 & 0.757 \\
Kolodziejczyk (3) & 4\textrm{-}12 & 0.722 & 0.630 & 0.722 & 0.722 & 0.838 & 0.838 \\
Klein (4) & 3\textrm{-}17 & 0.501 & 0.501 & 0.863 & 0.863 & 0.973 & 0.859 \\
Zeisel (7) & 3\textrm{-}20 & 0.518 & 0.518 & 0.648 & 0.648 & 0.855 & 0.855 \\
\midrule
\textit{Mean} & & \textbf{0.661} & \textbf{0.473} & \textbf{0.772} & \textbf{0.665} & \textbf{0.632} & \textbf{0.635} \\
\bottomrule
\end{tabular}
\end{table*}


\section{Mode against size}\label{app:modesize}
Table~\ref{tab:modesize} gives the three arms of the test in Section~\ref{sec:modesize}. Within a
seed the cells, the subsamples and the labels are shared. The pipelines differ in their
normalisation of the mutual information and in the cluster counts their grids reach.

\begin{table}[t]
\centering
\caption{Mode against size. One fixed pool of 2,500 PBMC14K cells, drawn afresh at each of three seeds, run through both of MICA's pipelines and the baseline under the same rule. The data, the size, the subsamples and the labels are identical across the three rows. The pipelines differ, and so do the cluster counts their grids reach: graph mode cannot make fewer than 4 or 5 clusters here, and Section~\ref{sec:modesize} shows that this, not the mode, accounts for the difference. Tied settings are those within 0.01 of the best. The span is the range of cluster counts among them, reported as the median over the three pools with the observed range. The last two columns say how often the two readings of the rule chose differently and by how much in ARI.}
\label{tab:modesize}
\footnotesize
\setlength{\tabcolsep}{4pt}
\resizebox{\columnwidth}{!}{%
\begin{tabular}{lcccc}
\toprule
Pipeline & Share tied & Span & Readings differ & Median gap \\
\midrule
MICA, small-data mode & 0.26 & 4 (2\textrm{-}5) & 3 of 3 & 0.319 \\
MICA, graph mode & 0.16 & 0 (0\textrm{-}1) & 0 of 3 & 0.000 \\
Leiden, as a control & 0.15 & 1 (1\textrm{-}1) & 1 of 3 & 0.000 \\
\bottomrule
\end{tabular}}
\end{table}


\section{The default and silhouette rules}\label{app:secondary}
The body concentrates on the stability rule's two readings. The default and the silhouette rule
carry choices of their own, and both favour MICA for reasons that have little to do with its
clustering.

\textbf{The shipped default.}
MICA's best result comes from its own shipped default, which returns four clusters whatever
the dataset: a mean ARI of 0.727, leading on seven of nine datasets. The median number of
cell types across these datasets is also four, and three have exactly four. Split by that,
the default scores 0.802 where the true count is 3 to 5 and 0.579 outside. We cannot show
how it was chosen, but tuning $k$ on the other eight datasets' labels picks 4 in 26 of 27 cases
(Appendix~\ref{app:transfer}), so the default is what tuning on this benchmark would select.
Its apparent strength is a property of this benchmark's composition rather than of the method,
and it would not carry over to tissue with a different number of types. The baseline's default is a resolution rather than a count, and it overclusters
badly here, reaching 21 clusters on Zeisel.

\textbf{The silhouette rule.} It carries a third such choice. Scoring the same partitions in the 50
principal components rather than in MICA's own embedding changes the selected ARI by 0.13
on average for MICA and 0.16 for the baseline, reaching 0.52, with only the space changing
(Appendix~\ref{app:repro}). On matched candidates, where both methods are offered the same
counts, MICA's mean holds at 0.593 while the baseline's falls from 0.564 to 0.478, since it then
picks 2 clusters on seven of the nine datasets.

\section{The transfer comparison}\label{app:transfer}
Table~\ref{tab:transfer} gives the transfer comparison of Section~\ref{sec:transfer} per
dataset. It uses the other datasets' labels to choose each setting, so it answers how far a
setting tuned on labelled benchmarks carries to a new one, not what a user could do without
labels. MICA's tuned $k$ is 4 on every held-out dataset except Yan at one seed, where it is 3.
The baseline's tuned resolution lies between 0.22 and 0.30 and gives 2 to 8 clusters.

\begin{table}[t]
\centering
\caption{The transfer comparison, which is not label-free. Each small dataset takes the setting with the best mean ARI over the other eight at the same seed (MICA: $k$, Leiden: a resolution), and PBMC14K takes the baseline's resolution best on the nine. MICA's graph mode has no setting that carries over from $k$. ARI with the resulting cluster count in parentheses, medians over three seeds.}
\label{tab:transfer}
\footnotesize
\setlength{\tabcolsep}{4pt}
\begin{tabular}{lcc}
\toprule
Dataset (types) & MICA & Leiden \\
\midrule
Yan (8) & 0.765\,(3\textrm{-}4) & 0.765\,(4) \\
Goolam (5) & 0.582\,(4) & 0.619\,(2) \\
Buettner (3) & 0.919\,(4) & 0.134\,(2) \\
Pollen (11) & 0.493\,(4) & 0.606\,(6\textrm{-}7) \\
Chung (4) & 0.655\,(4) & 0.510\,(6) \\
Usoskin (4) & 0.839\,(4) & 0.696\,(5) \\
Kolodziejczyk (3) & 0.834\,(4) & 0.696\,(5) \\
Klein (4) & 0.980\,(4) & 0.886\,(5) \\
Zeisel (7) & 0.478\,(4) & 0.855\,(8) \\
PBMC14K (7) & n/a & 0.786\,(6) \\
\midrule
\textit{Mean, nine small} & \textbf{0.727} & \textbf{0.641} \\
\bottomrule
\end{tabular}
\end{table}


\section{Adjusted mutual information}\label{app:ami}
Table~\ref{tab:ami} gives the AMI for the same chosen partitions as Table~\ref{tab:rulesall}.
The two measures agree on which method leads under every rule. The AMI sits above the ARI
whenever the chosen count is wrong, by 0.17 on average across MICA's stability picks with too
many clusters and 0.12 with too few, and it can be far more forgiving of a partition that merges
types: on Pollen MICA's 17-cluster partition scores 0.898 under ARI and 0.913 under AMI, while
its 2-cluster partition scores 0.087 and 0.318. Being the stricter of the two in both
directions is why we report the ARI in the main table.

\begin{table*}[t]
\centering
\caption{Adjusted mutual information for the same chosen partitions as Table~\ref{tab:rulesall}, on the same runs. The AMI and the ARI agree on which method leads under every rule. The AMI sits above the ARI whenever the chosen count is wrong, in either direction, and it is much more forgiving of a partition that merges cell types, such as MICA's 2-cluster pick on Pollen.}
\label{tab:ami}
\footnotesize
\setlength{\tabcolsep}{4pt}
\begin{tabular}{lcccccccc}
\toprule
 & \multicolumn{2}{c}{Default} & \multicolumn{2}{c}{Silhouette} & \multicolumn{2}{c}{Finest stable} & \multicolumn{2}{c}{Most stable} \\
\cmidrule(lr){2-3}\cmidrule(lr){4-5}\cmidrule(lr){6-7}\cmidrule(lr){8-9}
Dataset (types) & MICA & Leiden & MICA & Leiden & MICA & Leiden & MICA & Leiden \\
\midrule
Yan (8) & 0.804 & 0.846 & 0.817 & 0.804 & 0.721 & 0.846 & 0.652 & 0.804 \\
Goolam (5) & 0.749 & 0.587 & 0.679 & 0.652 & 0.771 & 0.584 & 0.679 & 0.652 \\
Buettner (3) & 0.884 & 0.272 & 0.590 & 0.304 & 0.904 & 0.287 & 0.590 & 0.283 \\
Pollen (11) & 0.731 & 0.917 & 0.943 & 0.924 & 0.913 & 0.890 & 0.318 & 0.890 \\
Chung (4) & 0.762 & 0.738 & 0.725 & 0.128 & 0.709 & 0.339 & 0.292 & 0.339 \\
Usoskin (4) & 0.798 & 0.624 & 0.508 & 0.755 & 0.821 & 0.709 & 0.821 & 0.747 \\
Kolodziejczyk (3) & 0.888 & 0.723 & 0.784 & 0.723 & 0.822 & 0.896 & 0.757 & 0.896 \\
Klein (4) & 0.963 & 0.700 & 0.664 & 0.890 & 0.664 & 0.956 & 0.664 & 0.890 \\
Zeisel (7) & 0.565 & 0.644 & 0.495 & 0.580 & 0.648 & 0.815 & 0.648 & 0.815 \\
PBMC14K (7) & 0.868 & 0.893 & n/a & n/a & 0.790 & 0.725 & 0.790 & 0.725 \\
\midrule
\textit{Mean, nine small} & \textbf{0.794} & \textbf{0.672} & \textbf{0.689} & \textbf{0.640} & \textbf{0.775} & \textbf{0.702} & \textbf{0.602} & \textbf{0.702} \\
\bottomrule
\end{tabular}
\end{table*}


\section{Reproducing the results}\label{app:repro}
\texttt{scripts/a3\_report.py} rebuilds the figure and the tables of the runs as first
specified from the saved result files and writes \texttt{results/a3\_report\_numbers.json}.
\texttt{scripts/a3\_matched.py} builds Tables~\ref{tab:matched}, \ref{tab:summary} and
\ref{tab:transfer} from the matched runs and writes \texttt{results/a3\_matched\_numbers.json}.
Between them the two files hold every aggregate figure quoted in the text, each traceable to the
run that produced it. Per-dataset ARI ranges across seeds are in
\texttt{results/a3\_seed\_variation.csv} and \texttt{results/a3\_matched\_seed\_variation.csv}.

\textbf{Software.} Two conda environments, both pinned in \texttt{envs/}. The analysis
environment has Python 3.11, Scanpy 1.11.5 and scikit-learn. The second reproduces the
authors' published environment, Python 3.9.2 with MICA 1.0.0 and scikit-learn 0.24.1, and
is used whenever MICA's own code runs. MICA is the released version at commit
\texttt{54fae81}.

\textbf{Running the experiment.} The stability rule for the nine small-data datasets is
\texttt{scripts/run\_rules\_small.sh}, which calls
\texttt{scripts/run\_rule\_stability.py} once per dataset at seed 0, with
\texttt{scripts/run\_rules\_repeats.sh} supplying seeds 1 and 2. PBMC14K needs MICA's
graph mode and runs through \texttt{scripts/run\_rule\_stability\_graph.py}, which shells
out to the legacy environment for the node2vec step. The mode against size comparison of
Section~\ref{sec:modesize} is \texttt{scripts/run\_rule\_mode\_size.py}, which draws the
pool once and calls both pipelines on it. The silhouette rule is
\texttt{scripts/run\_rule\_silhouette.py} and the shipped defaults are
\texttt{scripts/run\_rule\_default.py}. Each writes a folder under \texttt{results/}
containing \texttt{metrics.csv} and a \texttt{params.json} recording the settings, seed
and candidate grid of that run. The matched runs of Section~\ref{sec:matcheddesign} use the
same stability, mode and silhouette runners with \texttt{-{}-matched} on the command line,
called through \texttt{scripts/run\_rules\_matched.sh} once per part and seed, and write to
folders starting \texttt{rule\_matched\_}. Their grids are in
\texttt{src/scminer\_clust/grids.py}, fixed beforehand from the label-free sweeps of
\texttt{scripts/probe\_grid\_coverage.py} and \texttt{scripts/probe\_mica\_graph\_coverage.py}.

\textbf{Settings.} Subsamples: $B = 25$ at 80 percent of cells, redrawn per seed, shared
between the two methods within a seed. Candidate settings in the runs as first specified:
$k \in \{2,\dots,20\}$ for MICA
in small-data mode, resolutions 0.05 to 2.00 in steps of 0.05 for Leiden, and 0.1 to 4.0 in
steps of 0.1 for MICA in graph mode, which is MICA's own grid. Only settings whose
full-data partition yields 2 to 20 clusters enter the comparison. The matched runs sweep 300
log-spaced Leiden resolutions from 0.001 to 50 and, in MICA's graph mode, 0.1 to 2.0 in steps
of 0.01, then to 4.0 in steps of 0.1 and to 30 in steps of 0.5, and keep one setting per count
among the counts every method reaches. $k$-means uses 50
restarts. The bootstrap for the stability threshold is a percentile bootstrap over cells with
2,000 resamples, where chooseR specifies BCa with 25,000. The two differ only in the
lower bound they give, which is what sets the threshold. We first assumed this would not
matter and then tested it, with the result reported in Section~\ref{sec:results}: the
chosen cluster count is sensitive to shifts of this size, so the difference between the
two bootstraps could change the selection. This is the clearest unverified assumption
remaining in our implementation, and settling it needs a run against chooseR's own R
code.

\textbf{Scale of the computation.} Each dataset and seed runs $B+1$ clusterings per
candidate setting for each method, which is 494 for MICA over 19 values of $k$ and 468 to
1,040 for Leiden, whose grid loses settings that fall outside 2 to 20 clusters, and 364 to 494
for either method on matched candidates. The reported cost excludes the one-off mutual information
distance matrix, which is computed once per dataset, cached, and shared by every setting
and seed.


\textbf{The silhouette rule and the space it is scored in.} Each partition from both
methods was scored both in the 50 principal components and in MICA's multidimensional
scaling embedding. The selected ARI differs between the two spaces by 0.13 on average
for MICA and 0.16 for the baseline, with the largest differences on Goolam, 0.38 for
MICA, and Chung, 0.52 for the baseline. The partitions, the candidate settings and the
methods are identical across the two, so the difference is attributable to the space
alone. Per-dataset values are in the \texttt{rule\_silhouette} folders under
\texttt{results/}.

\section{Why subsampling the distance matrix is exact}\label{app:exact}
Section~\ref{sec:method} states that MICA's distance matrix restricted to a subset of
cells equals the matrix computed from that subset alone. The reason is that
Eq.~\ref{eq:mi} involves only the two cells being compared. Cell $i$ is discretised using
its own minimum and maximum across genes, so its bin codes $C_i$ do not change when other
cells are removed. The joint histogram of $C_i$ and $C_j$ is taken over the $g$ genes, which
are also unchanged. Hence every term in $I(C_i; C_j)$ and $H(C_i, C_j)$ is a function of
cells $i$ and $j$ alone. The same holds for the small-data form, Eq.~\ref{eq:misd}, whose
marginal entropies $H(C_i)$ and $H(C_j)$ are each computed from one cell's bin codes.

This fails for any method that standardises genes across cells, principal components
included, which is why the baseline is fully recomputed on every subsample while MICA's
distance is not. The claim is checked rather than assumed:
\texttt{tests/test\_subsample\_invariance.py} recomputes the distance on a random half of
a dataset and asserts exact equality with the corresponding submatrix, in both of MICA's
modes.

\section{Anomalies raised against the Assignment 2 tables}\label{app:anomalies}
A team member reviewing the result tables flagged six values as possible data handling
errors. Each was checked against the stored partitions and gene sets, and all six are
genuine results. We record them here because two are informative.

Identical scores in different gene settings do not indicate reused results. Assignment 2
also ran each method on the 2,000 most variable genes under three different selectors,
where this paper keeps all genes. On Goolam those three gene sets give hierarchical
clustering partitions that agree at ARI 1.000 with each other, so the repeated 0.910 is one
partition found three times. On
Pollen the three MICA partitions are different from one another, agreeing at ARI 0.91 to
0.95, yet each scores 0.942 against the labels, which is a coincidence of equal agreement
rather than a copied value.

The collapse of every method on Chung with the \texttt{seurat} gene set is also real. That
input is stored on a natural log scale, and after correcting the log base in Assignment 2
the \texttt{seurat} selector shares only 715 of its 2,000 genes with \texttt{cell\_ranger}.
The two gene sets are genuinely different, and one of them is poor for this dataset.
